% ============================================================
% Chapter 4: Experiments
% CT-MUSIQ Undergraduate Thesis
% ============================================================

\chapter{Experiments}
\label{chap:experiments}

\section{Overview}
\label{sec:exp_overview}

This chapter provides a comprehensive description of the experimental framework used to
evaluate CT-MUSIQ and to validate the contributions of each proposed architectural and
training component. Section~\ref{sec:exp_impl} details the software and hardware
environment. Section~\ref{sec:exp_hyperparams} documents the full set of hyperparameters
used for both CT-MUSIQ and the Agaldran Combo baseline. Section~\ref{sec:exp_training}
describes the main training runs in chronological order. Section~\ref{sec:exp_ablation}
presents the complete ablation study protocol and results. Section~\ref{sec:exp_eval}
specifies the evaluation procedure for the held-out test set. Section~\ref{sec:exp_repro}
discusses reproducibility provisions.


\section{Implementation Details}
\label{sec:exp_impl}

\subsection{Software Environment}

The complete CT-MUSIQ project is implemented in Python~3.12 running in a virtual
environment (\texttt{.venv}) created with \texttt{py -3.12 -m venv .venv}. The following
core dependencies are used:

\begin{itemize}
    \item \textbf{PyTorch~2.x}: Primary deep learning framework for model definition,
          training loop, automatic differentiation, and CUDA acceleration.
    \item \textbf{timm ($\geq$ 0.9.0)}: Model zoo library providing pretrained ViT-B/32
          weights (\texttt{vit\_base\_patch32\_224}) and Swin-T weights
          (\texttt{swin\_tiny\_patch4\_window7\_224}).
    \item \textbf{torchvision ($\geq$ 0.15)}: Used for pretrained ResNet-50 weights
          in the Agaldran Combo baseline.
    \item \textbf{PIL/Pillow}: Image loading and float-mode augmentation operations,
          including bicubic resizing and rotation without quantisation artefacts.
    \item \textbf{NumPy}: Array operations during preprocessing, augmentation, and
          metric computation.
    \item \textbf{SciPy}: Pearson, Spearman, and Kendall correlation coefficient
          computation for evaluation.
    \item \textbf{scikit-learn}: Isotonic regression calibration for the optional
          post-processing step.
    \item \textbf{pandas}: Training log CSV management for epoch-by-epoch logging.
    \item \textbf{tqdm}: Progress bar display during training and evaluation.
\end{itemize}

All hyperparameters and file paths are centralised in a single \texttt{config.py}
file, which serves as the sole source of truth for the entire experiment. Training,
evaluation, and ablation scripts import configuration constants from this file, ensuring
that no hyperparameter is ever hard-coded in a script and that all experiments are fully
reproducible by modifying only the configuration file.

\subsection{Hardware Specifications}

All experiments in this thesis are conducted on a single workstation with the following
hardware specifications:

\begin{itemize}
    \item \textbf{GPU}: NVIDIA GeForce RTX~3060 Laptop GPU (6~GB GDDR6 VRAM, Ampere
          architecture, 3840 CUDA cores, 30 TFLOPS fp16 Tensor Core performance).
    \item \textbf{CPU}: Intel Core i7 mobile processor.
    \item \textbf{RAM}: 16~GB DDR4 system memory.
    \item \textbf{Storage}: NVMe SSD for fast dataset loading (eliminates I/O bottleneck
          with \texttt{num\_workers = 0} on Windows).
    \item \textbf{Operating System}: Windows 11 (64-bit).
    \item \textbf{CUDA version}: 11.8 (compatible with PyTorch 2.x).
\end{itemize}

The RTX~3060 Laptop GPU's 6~GB VRAM budget is the principal hardware constraint that
motivated the mixed-precision training strategy, the two-scale pyramid (rather than
three-scale), and the batch size of 10. Training one epoch of CT-MUSIQ on the 700-image
training set requires approximately \textbf{31--40 seconds}. A complete 100-epoch
training run requires approximately \textbf{50--55 minutes}. The Agaldran Combo baseline
trains somewhat faster due to its single-resolution input pipeline.


\section{Hyperparameter Configuration}
\label{sec:exp_hyperparams}

Table~\ref{tab:hyperparams} summarises the complete hyperparameter configuration for
the primary CT-MUSIQ training run. All values are as specified in \texttt{config.py}.

\begin{table}[h]
    \centering
    \caption{Complete CT-MUSIQ hyperparameter configuration.}
    \label{tab:hyperparams}
    \begin{tabular}{lll}
        \toprule
        \textbf{Category} & \textbf{Hyperparameter} & \textbf{Value} \\
        \midrule
        Architecture & Embedding dimension $d$ & 768 \\
                     & Number of attention heads $H$ & 8 \\
                     & Number of Transformer layers $L$ & 12 \\
                     & FFN hidden dimension & 3072 \\
                     & Dropout rate & 0.05 \\
                     & Multi-scale sizes (pixels) & [224, 384] \\
                     & Patch size (pixels) & 32 \\
                     & Total tokens per image & 193 \\
                     & Max grid size $G_{\max}$ & 13 \\
        \midrule
        Training     & Batch size & 10 \\
                     & Maximum epochs & 100 \\
                     & Stage~1 epochs (frozen encoder) & 5 \\
                     & Stage~2 warm-up epochs & 3 \\
                     & Stage~1 learning rate & $3 \times 10^{-4}$ \\
                     & Stage~2 base learning rate & $5 \times 10^{-5}$ \\
                     & Minimum learning rate & $1 \times 10^{-6}$ \\
                     & Weight decay & 0.01 \\
                     & Gradient clipping norm & 1.0 \\
                     & EMA decay $\rho$ & 0.999 \\
                     & Early stopping patience & 20 \\
        \midrule
        Loss         & MSE loss weight & 1.0 \\
                     & KL loss weight $\lambda_{\text{KL}}$ & 0.10 \\
                     & Ranking loss weight $\lambda_{\text{rank}}$ & 0.50 \\
                     & Number of KL bins $K$ & 20 \\
                     & KL Gaussian $\sigma$ & 0.5 \\
                     & Ranking margin $m$ & 0.1 \\
        \midrule
        CT Window    & Window level (HU) & 40 \\
                     & Window width (HU) & 80 \\
        \midrule
        Dataset      & Training images & 700 (indices 0--699) \\
                     & Validation images & 200 (indices 700--899) \\
                     & Test images & 100 (indices 900--999) \\
        \midrule
        Reproducibility & Random seed & 42 \\
                        & cuDNN deterministic & True \\
        \bottomrule
    \end{tabular}
\end{table}

\subsection{Baseline Hyperparameters}

The Agaldran Combo baseline uses the same dataset split, loss function, and batch size
as CT-MUSIQ, but with the following differences:

\begin{itemize}
    \item \textbf{No two-stage training}: Trained end-to-end from epoch 1 with a constant
          learning rate of $3 \times 10^{-4}$, transitioning to cosine annealing from
          epoch~6 (same as CT-MUSIQ Stage~2 cosine phase).
    \item \textbf{Input resolution}: Single resolution ($224 \times 224$), not
          multi-scale patches.
    \item \textbf{EMA}: Applied identically with $\rho = 0.999$.
    \item \textbf{Early stopping patience}: 20 epochs, same as CT-MUSIQ.
\end{itemize}

The Agaldran Combo model has approximately $\sim$30~million parameters (Swin-T:
$\sim$28M, ResNet-50: $\sim$25M, shared prediction heads: $\sim$0.5M, noting that
the two backbones are trained jointly in a single model object).


\section{Training Runs}
\label{sec:exp_training}

\subsection{CT-MUSIQ Primary Training Run}

The primary CT-MUSIQ training run spans 100 epochs using the full two-stage schedule
and all hyperparameters in Table~\ref{tab:hyperparams}. The training log captures
per-epoch values of: training total loss, training MSE loss, training KL loss,
validation total loss, validation MSE, validation KL, PLCC, SROCC, KROCC, Aggregate,
and wall-clock time per epoch.

\textbf{Stage~1 convergence (epochs 1--5):} The validation aggregate rises from
approximately 2.51 at epoch~1 (where the encoder is frozen and only the randomly
initialised heads are trained) to approximately 2.59 by epoch~5. The training loss
falls steeply from $\sim$0.88 to $\sim$0.34 as the heads learn a crude quality mapping
from the frozen pretrained features.

\textbf{Stage~2 early phase (epochs 6--20):} Upon unfreezing the encoder, the
validation aggregate rises sharply to 2.71 at epoch~4 of Stage~2 (epoch~9 overall),
demonstrating that the pretrained encoder rapidly adapts its representations to the
CT quality assessment task. The training loss continues to decrease as the KL loss
becomes active and begins penalising scale inconsistencies.

\textbf{Stage~2 convergence phase (epochs 21--100):} The validation aggregate rises
steadily from $\sim$2.65 to a peak of approximately 2.76, with the best checkpoint
saved at the epoch achieving the maximum validation aggregate. The cosine learning rate
decay provides a smooth convergence trajectory without oscillation. Training MSE falls
from $\sim$0.27 to $\sim$0.12 over this phase.

\textbf{Best validation checkpoint}: The saved checkpoint achieves a validation
aggregate of $\sim$2.76, which corresponds to the model used for test set evaluation.

\subsection{Agaldran Combo Training Run}

The Agaldran Combo baseline is trained for 100 epochs under its simplified training
schedule. The training proceeds rapidly in the first 10 epochs as both backbones
adapt from ImageNet classification to CT quality regression. The ensemble of Swin-T
and ResNet-50 achieves strong performance due to the complementary feature representations
of window-based attention (Swin-T) and pyramid-pooled convolutional features (ResNet-50).

The best checkpoint for the Agaldran Combo model is saved at the validation aggregate
peak during training, and this checkpoint is used for final test set evaluation.


\section{Ablation Study}
\label{sec:exp_ablation}

\subsection{Protocol}

Six ablation configurations (A1--A6) are trained for 30 epochs each, using the same
random seed (42), training data, validation data, and preprocessing pipeline as the
primary CT-MUSIQ run. The shorter 30-epoch budget is chosen to enable evaluation of
all six configurations within the available compute budget while still reaching a
meaningful performance level for comparison. The validation aggregate at epoch~30 is
used as the comparison metric.

\subsection{Ablation Configuration Details}

\begin{description}
    \item[Configuration A1 (Single Scale, No KL):] Only the coarse scale ($224 \times
    224$, 49 patches) is used. The KL loss is disabled ($\lambda_{\text{KL}} = 0$),
    and the per-scale heads are absent. This configuration serves as the single-scale
    baseline that represents the minimum viable model, testing the benefit of multi-scale
    input against a single-scale ViT-based quality assessor.

    \item[Configuration A2 (Two Scales, No KL):] Both scales ($224 \times 224$ and
    $384 \times 384$, 193 total patches) are used with $\lambda_{\text{KL}} = 0$.
    Per-scale heads are present (to produce $\hat{q}_s$) but their disagreement is not
    penalised. This configuration isolates the contribution of the KL regularisation
    loss from the contribution of multi-scale tokenisation.

    \item[Configuration A3 (Two Scales, KL $\lambda = 0.05$):] The full architecture
    with two scales, but with a weaker KL weight. If A3 underperforms A4, the stronger
    $\lambda = 0.10$ provides necessary regularisation. If A3 matches A4, there is little
    sensitivity to the exact KL weight in this range.

    \item[Configuration A4 (Two Scales, KL $\lambda = 0.10$):] The proposed full
    CT-MUSIQ configuration. Serves as the reference point within the ablation.

    \item[Configuration A5 (Two Scales, KL $\lambda = 0.20$):] Tests whether doubling
    the KL weight over-regularises the model, potentially forcing the per-scale heads to
    agree at the expense of the global head's ability to learn the true quality mapping.

    \item[Configuration A6 (Three Scales, KL $\lambda = 0.10$):] Adds a third scale at
    $512 \times 512$ pixels (16$\times$16 = 256 additional patches, 449 total tokens).
    This requires updating the maximum grid size to $G_{\max} = 17$ and increases the
    sequence length from 193 to 449 tokens, substantially increasing VRAM usage. This
    tests whether native-resolution processing provides additional quality-relevant
    information beyond the two-scale representation.
\end{description}

\subsection{Ablation Results}

Table~\ref{tab:ablation_results} presents the validation aggregate at epoch~30 for all
six configurations. The results demonstrate several key findings:

\begin{table}[h]
    \centering
    \caption{Ablation study results (validation aggregate at epoch 30, 30-epoch training).}
    \label{tab:ablation_results}
    \begin{tabular}{llcccc}
        \toprule
        Config & Description & Scales & $\lambda_{\text{KL}}$ & Val.\ Aggregate \\
        \midrule
        A1 & Single scale, no KL   & [224]        & 0.00 & 2.249 \\
        A2 & Two scales, no KL     & [224, 384]   & 0.00 & 2.271 \\
        A3 & Two scales, KL~0.05   & [224, 384]   & 0.05 & 2.289 \\
        A4 & Two scales, KL~0.10   & [224, 384]   & 0.10 & 2.305 \\
        A5 & Two scales, KL~0.20   & [224, 384]   & 0.20 & 2.281 \\
        A6 & Three scales, KL~0.10 & [224, 384, 512] & 0.10 & 2.296 \\
        \bottomrule
    \end{tabular}
\end{table}

\textbf{Key observations:}

\begin{enumerate}
    \item \textbf{Multi-scale vs.\ single-scale (A1 vs.\ A2):} Adding the second scale
    (384~pixels) improves the 30-epoch aggregate from 2.249 to 2.271, a gain of +0.022.
    This confirms that the fine-scale patches ($144 \times$ $32 \times 32$ patches at
    384~px) provide complementary quality-relevant features not captured at the coarse
    scale alone.

    \item \textbf{KL regularisation benefit (A2 vs.\ A4):} Adding the scale-consistency
    KL loss at $\lambda = 0.10$ further improves aggregate from 2.271 to 2.305, a gain
    of +0.034 over the two-scale model without KL. This is the largest single improvement
    from any component tested in the ablation, confirming that scale-consistency
    regularisation is the most impactful innovation.

    \item \textbf{KL weight sensitivity (A3, A4, A5):} The optimal KL weight is
    $\lambda = 0.10$ (A4 = 2.305 $>$ A3 = 2.289 $>$ A5 = 2.281). A weight of 0.05 is
    insufficient to enforce strong cross-scale consistency, while a weight of 0.20
    begins to over-regularise by pushing the per-scale heads too strongly toward the
    global head at the expense of MSE accuracy. The performance difference between A3
    and A5 is relatively small (0.008), suggesting mild sensitivity to the exact KL
    weight in this range, but $\lambda = 0.10$ is robustly optimal.

    \item \textbf{Three scales vs.\ two scales (A4 vs.\ A6):} Adding the $512 \times 512$
    native-resolution scale slightly reduces the 30-epoch aggregate from 2.305 to 2.296.
    This marginal underperformance may be attributable to the increased sequence length
    (449 tokens) making training convergence slightly slower within the 30-epoch budget,
    rather than a true quality signal deficit. Full 100-epoch training of A6 would be
    needed to draw a definitive conclusion about the benefit of native-resolution
    processing.
\end{enumerate}

\subsection{Ablation Summary}

Based on the ablation study, the proposed CT-MUSIQ configuration (two scales at
$[224, 384]$ with $\lambda_{\text{KL}} = 0.10$) is confirmed as the best-performing
configuration under the 30-epoch ablation budget. The scale-consistency KL loss is the
most impactful single component, contributing approximately twice the performance gain of
the additional scale alone. These results validate the three core design principles of
CT-MUSIQ articulated in Chapter~\ref{chap:introduction}: domain-specific preprocessing
(implicit in all configurations), multi-scale global receptive field, and regularisation
through scale consistency.


\section{Evaluation Procedure}
\label{sec:exp_eval}

The evaluation procedure is identical for CT-MUSIQ and the Agaldran Combo baseline:

\begin{enumerate}
    \item Load the best saved checkpoint (EMA weights stored in the checkpoint file).
    \item Set the model to evaluation mode (\texttt{model.eval()}) and disable gradient
          computation (\texttt{torch.no\_grad()}).
    \item Iterate over the test set (\texttt{DataLoader} with \texttt{shuffle=False},
          \texttt{batch\_size=10}).
    \item For each batch, perform a forward pass and collect the global predicted scores
          $\hat{q}_i$ for all 100 test images.
    \item Compute PLCC, SROCC, KROCC using SciPy's \texttt{pearsonr}, \texttt{spearmanr},
          and \texttt{kendalltau} functions between the collected predictions and ground
          truth scores.
    \item Compute the aggregate score as $|\text{PLCC}| + |\text{SROCC}| + |\text{KROCC}|$.
    \item Write predictions and metrics to CSV files
          (\texttt{ct\_musiq\_predictions.csv} and \texttt{ct\_musiq\_results.csv}).
\end{enumerate}

Optional TTA (horizontal flip averaging) can be enabled via the \texttt{--tta} flag
in \texttt{evaluate.py}, and optional isotonic regression calibration can be enabled
via \texttt{--calibrate}. The primary reported results do not use TTA or calibration
to match the most direct comparison with the LDCTIQAC 2023 leaderboard.

The comparison table in \texttt{ct\_musiq\_results.csv} includes both the computed
CT-MUSIQ results and the published leaderboard results from Lee et al.\
\cite{lee2025ldctiqac}, allowing direct numerical comparison.


\section{Reproducibility}
\label{sec:exp_repro}

Full reproducibility of all experiments is ensured through the following measures:

\begin{itemize}
    \item \textbf{Fixed random seed}: All experiments use \texttt{SEED = 42}, seeding
          \texttt{torch.manual\_seed()}, \texttt{torch.cuda.manual\_seed\_all()},
          \texttt{numpy.random.seed()}, and \texttt{Python random.seed()} at the
          beginning of each script.
    \item \textbf{Deterministic cuDNN}: \texttt{torch.backends.cudnn.deterministic = True}
          ensures that CUDA operations produce identical results across runs.
    \item \textbf{Deterministic data split}: The 70/20/10 dataset split is
          deterministic and index-based (not randomly sampled), requiring no seed for
          splitting.
    \item \textbf{Centralised configuration}: All hyperparameters in \texttt{config.py}
          are version-controlled; no hyperparameter exists in any script other than as
          a \texttt{config.*} reference.
    \item \textbf{Checkpoint storage}: Best model checkpoints are stored in
          \texttt{results/ct\_musiq/ct\_musiq\_best.pth} and include the full model
          state dictionary, optimiser state, EMA shadow weights, and best aggregate
          value achieved.
\end{itemize}

Despite the above provisions, minor numerical differences ($< 0.001$ in correlation
metrics) may arise between runs due to non-deterministic GPU operations (particularly
in LayerNorm and attention on some CUDA versions). All reported results are from a single
canonical training run with the fixed seed configuration.

\section{Comparison Protocol Justification}
\label{sec:exp_comparison_protocol}

A critical methodological question in any empirical deep learning study is whether the
comparison between methods is fair. This section explicitly addresses the fairness of
the comparisons made in this thesis.

\subsection{Comparison with LDCTIQAC 2023 Leaderboard}

The LDCTIQAC 2023 leaderboard results published in Lee et al.\ \cite{lee2025ldctiqac}
represent the official scores of challenge participants on the test set (images
0900--0999, the same 100-image test set used in this thesis). The comparison is therefore
directly valid in terms of the test data used, the evaluation metrics applied, and the
test set size. However, several sources of potential unfairness should be acknowledged:

\textbf{Access to validation labels:} Challenge participants had access to only
the training labels (indices 0000--0699) and were evaluated on the test set without
access to test labels. In this thesis, the validation split (indices 0700--0899) is also
used for hyperparameter tuning and model selection, as explicitly described. This matches
the challenge setup: validation labels were available to challenge participants for
tuning purposes.

\textbf{Computational resources:} Challenge participants may have had access to
significantly more powerful hardware (multi-GPU clusters, large-memory servers) than the
single consumer GPU used in this thesis. Methods that require multiple GPUs, large batch
sizes, or expensive ensembles at test time may achieve higher performance than what is
reproducible on consumer hardware. CT-MUSIQ is specifically designed to operate within
6~GB VRAM, which is a considerably tighter budget than typical challenge participants
would use.

\textbf{Development time:} Challenge submissions represent the best efforts of their
respective teams over the challenge duration, which typically includes extensive
hyperparameter search and model selection. CT-MUSIQ is developed within the scope of an
undergraduate thesis with a constrained timeline. The competitive performance achieved
despite this constraint is therefore notable.

\subsection{Within-Paper Comparison: CT-MUSIQ vs.\ Agaldran Combo}

The within-paper comparison between CT-MUSIQ and the Agaldran Combo baseline is
explicitly designed to be fair: both models are trained on the same data split, evaluated
on the same test set, using the same evaluation script and metrics. The primary
methodological difference is the model architecture (multi-scale patch Transformer vs.\
full-image ensemble), with training details adapted appropriately for each architecture
(two-stage schedule for CT-MUSIQ, simpler schedule for Agaldran Combo). The loss
function is identical for both models.

The only potential source of unfairness is that CT-MUSIQ uses the two-stage training
schedule, which involves an additional design choice (the Stage~1 epoch count and
learning rate) that was tuned informally based on observed training dynamics. The
Agaldran Combo's simpler training schedule requires fewer design choices and may be
more robustly configured. A fully controlled comparison would use identical training
schedules for both methods, which is left for future work.

\section{Additional Experimental Notes}

\subsection{Windows DataLoader Compatibility}

PyTorch's multiprocessing-based \texttt{DataLoader} is incompatible with Windows when
\texttt{num\_workers > 0} in some configurations due to the Windows process forking
model. To ensure compatibility, \texttt{num\_workers = 0} is used throughout, which
serialises data loading. This does not create a significant performance bottleneck
because the primary compute cost is in the Transformer forward and backward passes, and
the dataset fits entirely in system RAM (1,000 TIFF files $\times$ $\sim$1~MB each
$\approx$ 1~GB).

\subsection{Memory Profiling}

During the full CT-MUSIQ training run (batch size 10, two scales, AMP enabled), the
peak GPU memory allocation is approximately 5.2--5.6~GB, leaving a small safety margin
below the 6~GB VRAM limit. The memory breakdown is approximately: model parameters
in fp16~$\approx$~350~MB, activation tensors in fp16 for 193 tokens × batch
10~$\approx$~2~GB, gradient tensors in fp32~$\approx$~700~MB, optimiser states
(AdamW: two momentum buffers)~$\approx$~1.4~GB.

\subsection{Training with Resume}

The training script supports checkpoint resumption via the \texttt{--resume} flag,
which loads the latest checkpoint and continues training from the saved epoch. This
was used in practice to restart training after system interruptions, with no noticeable
effect on final performance due to the deterministic training configuration.
